\section{Executive function in ADHD}\label{sec:adhd}
ADHD is heterogeneous: executive difficulties are common but not universal. A meta-analysis of 83 studies found that groups with ADHD underperform controls on every executive task class examined, most strongly on response inhibition, vigilance, working memory and planning, with medium effect sizes ($0.46$--$0.69$) \citep{willcutt2005ef}. The same authors caution that executive weaknesses are neither necessary nor sufficient to cause all cases of ADHD, and fewer than half of the (mostly child) samples show impairment on any single executive test \citep{willcutt2005ef}. Measured more carefully, impairment is more prevalent: in a sample of 136 children, 89\% of those with ADHD were impaired on at least one executive function (62\% working memory, 27\% inhibitory control, 38\% set shifting) \citep{kofler2019heterogeneity}.

Anxiety is common in adult ADHD: a meta-analysis of 58 studies ($N=18{,}821$) found higher anxiety scores in adults with ADHD than without ($g=0.77$; $g=0.95$ against non-clinical and $g=0.44$ against clinical controls) \citep{hoskins2026anxiety}. Whether anxiety adds executive load in ADHD is less clear: in a sample of 339 children, trait anxiety did not uniquely predict working memory or inhibitory control, nor moderate the ADHD effect \citep{marsh2024anxiety}. The stacking assumption (low baseline executive capacity, lowered further by anxiety) is therefore plausible but, in children at least, not supported by the one careful test we found.
